\begin{ex}
Tính: $1 + 2 + 3 + 4$.
\shortans{10}
\loigiai{
Cộng lần lượt: $1 + 2 = 3$; $3 + 3 = 6$; $6 + 4 = 10$.
\textbf{Đáp số:} $10$.
}
\end{ex}

\begin{ex}
Tính: $10 - 1 - 2 - 3$.
\shortans{4}
\loigiai{
Trừ lần lượt: $10 - 1 = 9$; $9 - 2 = 7$; $7 - 3 = 4$.
\textbf{Đáp số:} $4$.
}
\end{ex}

\begin{ex}
Tìm số thích hợp điền vào ô trống: $2 + \dots + 1 = 9$.
\shortans{6}
\loigiai{
Ta có: $2 + 1 = 3$.\nSố cần điền là: $9 - 3 = 6$.
\textbf{Đáp số:} $6$.
}
\end{ex}

\begin{ex}
Tìm số thích hợp điền vào ô trống: $10 - 3 - \dots = 2$.
\shortans{5}
\loigiai{
Ta có: $10 - 3 = 7$.\nSố cần điền là: $7 - 2 = 5$.
\textbf{Đáp số:} $5$.
}
\end{ex}

\begin{ex}
Tìm $x$ biết: $x + x = 8$. Giá trị của $x$ là bao nhiêu?
\shortans{4}
\loigiai{
Hai số giống nhau có tổng bằng 8 là 4 (vì $4 + 4 = 8$). Vậy $x = 4$.
\textbf{Đáp số:} $4$.
}
\end{ex}

\begin{ex}
Tìm $y$ biết: $y + y = 10$. Giá trị của $y$ là bao nhiêu?
\shortans{5}
\loigiai{
Hai số giống nhau có tổng bằng 10 là 5 (vì $5 + 5 = 10$). Vậy $y = 5$.
\textbf{Đáp số:} $5$.
}
\end{ex}

\begin{ex}
Hùng có 6 viên bi, An có ít hơn Hùng 2 viên bi. Hỏi cả hai bạn có bao nhiêu viên bi?
\shortans{10}
\loigiai{
Số bi của An là: $6 - 2 = 4$ (viên bi).\nCả hai bạn có tất cả: $6 + 4 = 10$ (viên bi).
\textbf{Đáp số:} $10$.
}
\end{ex}

\begin{ex}
Mẹ có 10 cái bánh. Mẹ cho anh 3 cái, cho em 4 cái. Hỏi mẹ còn lại bao nhiêu cái bánh?
\shortans{3}
\loigiai{
Số bánh mẹ đã cho hai anh em là: $3 + 4 = 7$ (cái bánh).\nSố bánh mẹ còn lại là: $10 - 7 = 3$ (cái bánh).
\textbf{Đáp số:} $3$.
}
\end{ex}

\begin{ex}
Một cây bưởi có 8 quả. Bác nông dân hái đi 5 quả, sau đó trên cây lại lớn thêm 4 quả mới. Hỏi lúc này cây bưởi có bao nhiêu quả?
\shortans{7}
\loigiai{
Số quả bưởi trên cây lúc này là: $8 - 5 + 4 = 7$ (quả).
\textbf{Đáp số:} $7$.
}
\end{ex}

\begin{ex}
Trong một phòng học có 10 cái bàn. Có 3 cái bàn bị hỏng đưa đi sửa, sau đó mang về thêm 2 cái bàn mới. Hỏi trong phòng có bao nhiêu cái bàn?
\shortans{9}
\loigiai{
Số cái bàn có trong phòng là: $10 - 3 + 2 = 9$ (cái bàn).
\textbf{Đáp số:} $9$.
}
\end{ex}

\begin{ex}
Tìm số tự nhiên $a$ sao cho: $3 + a = 10 - 1$.
\shortans{6}
\loigiai{
Vế phải: $10 - 1 = 9$.\nTa có: $3 + a = 9 \Rightarrow a = 9 - 3 = 6$.
\textbf{Đáp số:} $6$.
}
\end{ex}

\begin{ex}
Tìm số tự nhiên $b$ sao cho: $8 - b = 2 + 2$.
\shortans{4}
\loigiai{
Vế phải: $2 + 2 = 4$.\nTa có: $8 - b = 4 \Rightarrow b = 8 - 4 = 4$.
\textbf{Đáp số:} $4$.
}
\end{ex}

\begin{ex}
Có 9 bạn học sinh đi dạo. Có 4 bạn rẽ sang trái, 3 bạn rẽ sang phải. Hỏi có bao nhiêu bạn đi thẳng?
\shortans{2}
\loigiai{
Tổng số bạn rẽ sang trái và phải là: $4 + 3 = 7$ (bạn).\nSố bạn đi thẳng là: $9 - 7 = 2$ (bạn).
\textbf{Đáp số:} $2$.
}
\end{ex}

\begin{ex}
Một sợi ruy băng dài 10 cm. Bạn Linh cắt lấy một đoạn 3 cm làm nơ, cắt tiếp một đoạn 4 cm gói quà. Hỏi sợi ruy băng còn lại dài bao nhiêu xăng-ti-mét?
\shortans{3}
\loigiai{
Tổng độ dài đã cắt là: $3 + 4 = 7$ (cm).\nĐộ dài sợi ruy băng còn lại là: $10 - 7 = 3$ (cm).
\textbf{Đáp số:} $3$.
}
\end{ex}

\begin{ex}
Tìm một số biết rằng lấy số đó cộng với chính nó rồi cộng thêm 2 thì được 10.
\shortans{4}
\loigiai{
Lấy 10 trừ đi 2 được 8: số đó cộng với chính nó bằng 8.\nSố cần tìm là $4$ (vì $4 + 4 = 8$).
\textbf{Đáp số:} $4$.
}
\end{ex}

\begin{ex}
Hai số có tổng bằng 10 và hiệu bằng 4. Hỏi số lớn hơn là số mấy?
\shortans{7}
\loigiai{
Hai số có tổng bằng 10 và hiệu bằng 4 là 7 và 3 (vì $7 + 3 = 10$ và $7 - 3 = 4$).\nSố lớn hơn là số $7$.
\textbf{Đáp số:} $7$.
}
\end{ex}

\begin{ex}
Hai số có tổng bằng 8 và hiệu bằng 2. Hỏi số bé hơn là số mấy?
\shortans{3}
\loigiai{
Hai số có tổng bằng 8 và hiệu bằng 2 là 5 và 3 (vì $5 + 3 = 8$ và $5 - 3 = 2$).\nSố bé hơn là số $3$.
\textbf{Đáp số:} $3$.
}
\end{ex}

\begin{ex}
Điền số thích hợp: $5 + 3 - \dots = 4$.
\shortans{4}
\loigiai{
Ta có: $5 + 3 = 8$.\nSố cần điền là: $8 - 4 = 4$.
\textbf{Đáp số:} $4$.
}
\end{ex}

\begin{ex}
Điền số thích hợp: $10 - 4 + \dots = 9$.
\shortans{3}
\loigiai{
Ta có: $10 - 4 = 6$.\nSố cần điền là: $9 - 6 = 3$.
\textbf{Đáp số:} $3$.
}
\end{ex}

\begin{ex}
Trên sân có 4 con gà và 3 con vịt. Bác nông dân mua thêm 2 con gà nữa. Hỏi trên sân có bao nhiêu con gà cả thảy?
\shortans{6}
\loigiai{
Số con gà trên sân có tất cả là: $4 + 2 = 6$ (con gà).
\textbf{Đáp số:} $6$.
}
\end{ex}

\begin{ex}
Tìm số tự nhiên lớn nhất có một chữ số thỏa mãn: khi cộng số đó với 1 thì kết quả vẫn bé hơn 10.
\shortans{8}
\loigiai{
Số đó cộng với 1 bé hơn 10 nghĩa là kết quả lớn nhất bằng 9.\nSố tự nhiên lớn nhất đó là: $9 - 1 = 8$.
\textbf{Đáp số:} $8$.
}
\end{ex}

\begin{ex}
Tìm số tự nhiên bé nhất khác 0 thỏa mãn: lấy 10 trừ đi số đó được kết quả lớn hơn 8.
\shortans{1}
\loigiai{
Ta có: $10 - 1 = 9 > 8$. Vậy số bé nhất khác 0 thỏa mãn là số $1$.
\textbf{Đáp số:} $1$.
}
\end{ex}

\begin{ex}
Mai có 7 cái bút, Nam có 3 cái bút. Mai phải cho Nam mấy cái bút để hai bạn có số bút bằng nhau?
\shortans{2}
\loigiai{
Mai nhiều hơn Nam số bút là: $7 - 3 = 4$ (cái bút).\nMai cần cho Nam số bút là: $4 : 2 = 2$ (cái bút).\nKiểm tra: Khi đó mỗi bạn có 5 cái bút.
\textbf{Đáp số:} $2$.
}
\end{ex}

\begin{ex}
Một đàn chim có 6 con bồ câu trắng và 4 con bồ câu nâu. Hỏi số bồ câu trắng nhiều hơn số bồ câu nâu bao nhiêu con?
\shortans{2}
\loigiai{
Số bồ câu trắng nhiều hơn số bồ câu nâu là: $6 - 4 = 2$ (con).
\textbf{Đáp số:} $2$.
}
\end{ex}

\begin{ex}
Tính giá trị của biểu thức: $10 - 2 + 1 - 3$.
\shortans{6}
\loigiai{
Thực hiện lần lượt: $10 - 2 = 8$; $8 + 1 = 9$; $9 - 3 = 6$.
\textbf{Đáp số:} $6$.
}
\end{ex}
